%==============================================================================
% Section 8: Torsion/Nieh-Yan sector robustness
%==============================================================================
\section{Torsion/Nieh--Yan sector robustness}
\label{sec:torsion}

Although the action contains the NY term and torsion, the theorem and robustness
checks so far have not referred to torsion.  In this section we examine the
effect of the torsion sector on the preceding results.  The claim is confined to
a single point: \emph{the torsion sector affects none of the theorem hierarchy of
this paper (quadratic / $k>0$ spectral completion / finite deformation).}  The
physics of torsion itself (boundary charge, localisation, etc.) is outside the
scope of this paper.

\subsection{Known/new boundary}
\label{sec:torsion:boundary}

That the NY term is topological in the canonical analysis and does not change the
classical dynamics is known~\cite{ref33,ref34,ref35}.  The content of this
section is not a restatement of that generality but an \emph{exact demonstration
in the reduced system}: (i) the reduced action is exactly at most quadratic in
the torsion ansatz scalars; (ii) under field promotion no derivative coupling
arises in the inspected criterion; and (iii) the finite-deformation
classification is invariant.

\subsection{(i) Exact auxiliary shell}
\label{sec:torsion:shell}

The reduced torsion Lagrangian of the NY-MAIN route can be organised as an
auxiliary shell in the form
\begin{equation}
  \Laux
  = a_{\eta}\,\eta^{2} + a_{V}\,V^{2} + b_{V}\,\dot{\vol}\,V
    + \text{(boundary total-derivative)},
\end{equation}
(the coefficients $a_{\eta},a_{V},b_{V}$ are topology/mode dependent but
algebraic in $\eta,V$).  By a polynomial-degree verification (symbolic up to
degree 4, residual zero), it is \emph{exactly at most quadratic} in $(\eta,V)$.

\paragraph{Algebraic proof (shell condition).}
The Euler--Lagrange equations are
\begin{equation}
  \frac{\partial\Laux}{\partial\eta}=2a_{\eta}\,\eta=0,\qquad
  \frac{\partial\Laux}{\partial V}=2a_{V}\,V + b_{V}\,\dot{\vol}=0.
\end{equation}
If $a_{\eta}\ne0$ then $\eta^{\ast}=0$.  The linear term
$b_{V}\,\dot{\vol}\,V$ in $V$ can be separated as the boundary current
\begin{equation}
  2\,\frac{\dd}{\dd t}\!\left(\vol\cdot V\right),
\end{equation}
and the bulk equation gives $V^{\ast}=0$.  Hence
\begin{equation}
  \eta^{\ast}=0,\qquad V^{\ast}=0
\end{equation}
is an \emph{exact auxiliary shell}, not an approximation
(Appendix~\ref{app:torsion}).

\subsection{(ii) Field promotion does not generate kinetic terms}
\label{sec:torsion:promotion}

Decompose the connection as $\Gamma=\Gamma_{\mathrm{LC}}+K(\eta,V)$, with
\begin{equation}
  R[\Gamma]=\dd\Gamma+\Gamma\wedge\Gamma,\qquad
  K=K(\eta,V)\ \text{(algebraic in the $\AX/\VT/\MX$ ansatz scalars)}.
\end{equation}
If $K$ is algebraic in $\eta,V$ then $dK$ is first order in the field
derivatives, and the quadratic derivative couplings
$(\dot\eta^{2},\dot V^{2},(\nabla\eta)^{2},(\nabla V)^{2})$ in the scalar
contraction of $\Delta R$ cancel structurally (a native-action-gate statement;
not a full torsion field theory).

Promoting the scalars to fields $\eta(x),V(x)$, we inspected the above over all
cells of four topologies $\times$ $\AX/\VT/\MX$.  The results:

\begin{itemize}
  \item No quadratic derivative coupling appears at all.  Even a linear spatial
        gradient is absent.
  \item What appears is only a boundary-type time-linear current: $2\,\vol\cdot\dot V$
        on the EC side and $6\,\vol^{2/3}\,\dot\eta$ on the NY side.
  \item At the level of the Riemann components individual derivative terms fire,
        but they cancel under the scalar contraction, as confirmed by a
        component computation.
  \item The classification procedure is validated by positive-detector
        controls in the native-action gate that the witness checks before
        reporting the cells above: a free-scalar toy run through the same
        classifier recovers the expected positive time-kinetic and negative
        spatial-gradient coefficients (alongside ghost, wrong-sign, and
        linear-current toys), so a genuine kinetic term would be detected
        (\code{script/scripts/paper07/torsion_auxiliary_field_promotion_witness.py};
        no claim about a full 24-component torsion field theory).
\end{itemize}

\subsection{(iii) Finite fate is unchanged}
\label{sec:torsion:finite}

The modulus/wall split of Section~\ref{sec:finite} is invariant under all of the
torsion modes EH/AX/VT/MX (a corollary of $\Laux|_{\mathrm{shell}}=0$ on the
auxiliary shell).

\subsection{Conclusion and scope}
\label{sec:torsion:conclusion}

\begin{quote}
Within the NY-MAIN route and the $\AX/\VT/\MX$ torsion ansatz scalars, the
reduced torsion sector remains an exact auxiliary, and its promotion to fields
$\eta(x),V(x)$ generates no quadratic kinetic/gradient term in the inspected
native-action criterion.  Hence torsion/NY affects none of the DBR rank-defect
kernel (Section~\ref{sec:rankdefect}), the robustness from the $k>0$ spectral
completion (Section~\ref{sec:spectral}), or the finite homogeneous modulus/wall
split (Section~\ref{sec:finite}).
\end{quote}

\paragraph{Scope (explicit).}
The subject of this section is the $\AX/\VT/\MX$ ansatz scalars, not a
$(\partial T)^{2}$-absence theorem for a full 24-component torsion field theory.
The exact-degree-2 / exact-shell claim is restricted to the NY-MAIN route.  A
time-linear boundary current is not called a kinetic/propagating mode.  The
absence of NY bulk involvement is not claimed outside this scope.

\begin{table}[H]
\centering
\begingroup
\renewcommand{\arraystretch}{1.4}
\begin{tabularx}{\linewidth}{lLLl}
\toprule
Layer & Check & Result & Scope \\
\midrule
Quadratic & Polynomial-degree verification ($\eta,V$) &
exact degree-2; $\eta^{\ast}=V^{\ast}=0$ exact shell & NY-MAIN route \\
Finite & Invariance of the modulus/wall split &
invariant (corollary) & finite Palatini branch \\
Field-promoted & Derivative-coupling classification (4 topo.\ $\times$ AX/VT/MX) &
kinetic/gradient absent; boundary current only & inspected native-action
criterion \\
\bottomrule
\end{tabularx}
\endgroup
\caption{Hierarchy of the torsion robustness.}
\label{tab:torsion}
\end{table}
